%=============================================================================
% KNOWN LIMITATIONS — UIDT v3.9
% \label{sec:limitations}
%
% Consolidated from CLAIMS.json (UIDT-C-034, C-041, C-046, C-050, C-052),
% CONSTANTS.md v3.9.5, and audit history.
% TKT-2026-05-12-manuscript-elite-revision Step 2/4
%=============================================================================
\section{Known Limitations and Open Questions}
\label{sec:limitations}

\noindent
Transparency is a core \textsc{uidt} principle. The following limitations
are not incidental caveats: they demarcate precisely where the framework
makes contact with empirical data and where it does not. Each limitation
carries an evidence category and, where applicable, a kill-switch criterion.

\subsection{L1 --- Geometric Suppression Factor}
\label{lim:L1}
\begin{limitation}
The holographic suppression of the cosmological constant involves a
geometric factor of order $10^{10}$. While the functional form is derived
within the \textsc{uidt} information-density framework, the precise
numerical value of this factor is calibrated, not predicted from first
principles. Classification: \evC{}.
\end{limitation}

\subsection{L2 --- Electron Mass Residual}
\label{lim:L2}
\begin{limitation}
The \textsc{uidt} cascade model produces a mass hierarchy consistent with
the observed quark and lepton spectrum at the order-of-magnitude level.
However, the predicted electron mass residual deviates from the PDG value
by approximately $0.3\%$. This residual is not yet explained within the
formal framework and constitutes a falsification candidate.
Classification: \evD{}.
\end{limitation}

\subsection{L3 --- Cosmological Calibration Sensitivity}
\label{lim:L3}
\begin{limitation}
The cosmological predictions of \textsc{uidt} (dark-energy equation of
state $w_0 \approx \wZero$, Hubble constant $H_0 \approx \HubbleUIDT\kms$)
are calibrated to DESI~DR2 data. The Hubble tension ($\sim 8\sigma$) and
$S_8$ tension ($\sim 3\sigma$) are \emph{not} declared resolved;
they are treated as calibration targets under active investigation.
Classification: \evC{}.
\end{limitation}

\subsection{L4 --- Gamma Parameter Origin}
\label{lim:L4}
\begin{limitation}
The kinetic invariant $\gamma_{\mathrm{kin}} \approx \GammaKin$ is
established phenomenologically and is numerically stable across the
\textsc{uidt} equation system. A derivation of $\gamma$ from first
principles (SU(3) gauge structure alone) is not available. The discrepancy
$\delta\gamma = \DeltaGammaVal$ between $\gamma$ and the bare SU(3)
candidate $\gamma_{\mathrm{bare}} \approx \GammaBare$ is documented in
CLAIMS.json (UIDT-C-041) and remains an open question. Classification: \evAminus{}.
\end{limitation}

\subsection{L5 --- Cascade Length $N = 99$}
\label{lim:L5}
\begin{limitation}
The mass-hierarchy cascade uses $N = 99$ iteration steps as an
empirically calibrated parameter. An alternative cascade architecture
with $N \approx 94.05$ has been proposed in internal audit
(UIDT-C-046, C-050) and is currently classified \evE{}. Neither value
is derived from the SU(3) gauge structure; both remain calibration choices
until a geometric counting argument is established.
\end{limitation}

\subsection{L6 --- Torsion Binding Energy Independence}
\label{lim:L6}
\begin{limitation}
The Torsion Binding Energy $E_T = \ETorsionfull$ is derived within the
\textsc{uidt} geometric framework and calibrated to cosmological data.
No independent experimental confirmation exists at the time of writing.
The value is consistent with the cosmological constant calibration but
cannot be considered experimentally verified. Classification: \evC{}.
\end{limitation}

\subsection{L7 --- Clay Millennium Prize Boundary}
\label{lim:clay}
\begin{limitation}
The \textsc{uidt} derivation of $\Delta^*$ constitutes a formal proof
\emph{within the \textsc{uidt} axiom system}. It does not claim to satisfy
the requirements of the Clay Mathematics Institute Millennium Prize
formulation, which demands a proof within the rigorous mathematical
framework of quantum field theory as specified in the Clay problem
statement. These are distinct claims and must not be conflated.
\end{limitation}

\subsection{L8 --- Born Rule Derivation (Research Programme)}
\label{lim:L8}
\begin{limitation}
Section~\ref{sec:born} presents a working formulation for deriving Born
rule statistics from the \textsc{uidt} information grid. As of v3.9, this
consists of a theorem statement and three identified technical tools
(T1--T3). The proof is \emph{not} complete. The claim that \textsc{uidt}
``resolves'' the measurement problem is explicitly false; the framework
presents a candidate mechanism under active construction.
Classification: \evE{}.
\end{limitation}
